\section{자료와 표본}

자료는 한국마사회 공개 경주자료에서 수집한 승식별 최종 배당률이다. 삼쌍승식이
도입된 2016년 6월 10일부터 2025년 12월 31일까지를 기본 범위로 한다. 코로나19
기간의 시장 운영이 평시와 크게 달랐던 2020년과 2021년은 제외한다. 이 범위의
파싱된 후보 자료는 2016--2019년과 2022--2025년의 19,301개 경주다. 원자료에
9,999를 넘는 값이 확인된 2018년 7월 1일의 17개 경주를 사전 규칙에 따라
제외하면 19,284개 경주가 남는다. 최종 분석표본과 조건별 제외 경주 수는 재현
스크립트가 생성한 \texttt{outputs/sample\_flow.csv}에서 확인할 수 있다.

분석 단위는 경주--승식--조합이다. 각 경주에서 유효 출전마 목록을 먼저 고정한
뒤, 이 목록으로 가능한 조합의 완전한 지지집합을 생성한다. 취소마, 판매중지
조합, 중복 레코드, 0 이하이거나 유한하지 않은 배당률을 순서대로 처리한다.
삼쌍승식의 관측 조합 수가 $n_r(n_r-1)(n_r-2)$와 일치하지 않는 경주는 결측
원인을 기록하고 주 분석에서 제외한다. 단순히 존재하는 행만 다시 정규화하면
결측 상태의 질량을 나머지 상태에 잘못 배분할 수 있기 때문이다.

주분석은 완전한 가격 지지집합을 요구하지만 행동모형에는 실현 착순도 필요하다.
행동모형 표본에서는 1--3위 착순과 유효 출전마의 대응을 추가로 검증한다.
조건부 순위확률을 추정하는 자료와 구조모형의 가격예측을 평가하는 자료는 연도를
fold로 나누어 교차적합한다. 따라서 같은 경주의 실현 착순으로 그 경주의
객관확률을 추정한 뒤 같은 가격에 맞추는 순환을 허용하지 않는다.

자료 검증은 다음 네 층으로 나눈다.

\begin{enumerate}
  \item \textbf{키 검증}: 날짜, 경마장, 경주번호, 승식, 조합이 유일한지 확인한다.
  \item \textbf{지지집합 검증}: 유효 출전마와 승식별 가능한 조합이 일치하는지 확인한다.
  \item \textbf{가격 검증}: 배당률의 단위, 최솟값, 게시 상한과 소수 첫째 자리 표기 구조를 확인한다.
  \item \textbf{재현 검증}: 원자료 파싱부터 논문의 표·그림까지 하나의 명령으로 다시 만든다.
\end{enumerate}

감사 결과, 사전 날짜 규칙을 통과한 19,284개 경주에서는 단승·쌍승·복승·
삼복승·삼쌍승 모두 가능한 조합 수와 관측 조합 수가 일치했다. 중복 조합,
유효 출전마 밖의 조합, 0 이하 또는 유한하지 않은 배당과 시장상태 불일치도
발견되지 않았다. 따라서 지지집합 결측에 따른 추가 제외는 없다. 이 판정은
\texttt{python -m analysis.data\_audit --strict}가 전체 조합키를 경주별로
검사한 결과다. 또한 \texttt{python -m analysis.display\_precision\_audit}는
상한이 아닌 게시 배당을 배당 크기 구간별로 나누어 모두 $0.1$ 격자 위에 놓이는지
별도로 검사하며, 한 행이라도 이 격자를 벗어나면 CI를 실패시킨다.

% Generated by python -m analysis.data_audit --strict; do not edit.
\begin{table}[htbp]
  \centering
  \caption{승식별 가격자료 감사와 분석표본}
  \label{tab:data-audit}
  \small
  \begin{tabular}{lrrrr}
    \hline
    승식(역할) & 완전·유효 & 승식 상한 & clean 점표본 & 상한 구간표본 \\
    \hline
    단승 (win) & 19,301 & 0 & 3,338 & 15,963 \\
    쌍승 (exacta) & 19,301 & 581 & 3,338 & 15,963 \\
    복승 (quinella) & 19,301 & 15 & 3,338 & 15,963 \\
    삼복승 (trio) & 19,301 & 3,434 & 3,338 & 15,963 \\
    삼쌍승 (원천) & 19,301 & 15,963 & 3,338 & 15,963 \\
    \hline
  \end{tabular}
  \begin{minipage}{0.96\linewidth}
    \footnotesize\emph{주:} 경주 수를 보고한다. 모든 행은 사전 제외일을
    뺀 19,284경주를 모집단으로 한다. 목표 승식 행은 이 중 삼쌍승과 해당
    승식의 완전한 지지집합 및 양수·유한 배당을 만족하는 경주 수다.
    clean 점표본은 두 승식 모두 상한이
    없는 경주다. 네 목표 승식의 clean 표본은 동일한 3,321개 race\_id
    집합이며, 삼쌍승 원천행에서 보듯 이는 19,284-15,963이다.
  \end{minipage}
\end{table}


게시 배당률은 풀 금액 그 자체가 아니다. 따라서 식
\eqref{eq:trifecta-state-price}은 투자자의 정확한 주관확률을 복원한다기보다
공제와 단수처리가 섞인 가격지표를 동일 합계로 맞춘다. 가능한 경우 원별 총매출과
공식 공제율을 이용한 대체 정규화를 보고하고, 불가능한 경우 그 한계를 결과의
해석 범위에 명시한다. 배당상한에 걸린 조합은 상한값을 정확한 가격으로 간주하지
않고 검열구간으로 처리한다.

상한이 아닌 게시 배당에 사용하는 $\pm0.05$ 구간은 단순한 편의적 가정이 아니다.
「한국마사회법 시행령」 제11조 제1항은 제10조에 따른 환급금에 10원 미만의
단수가 있을 때 5원 미만은 계산하지 않고 5원 이상은 10원으로 계산하도록
규정한다. 단위마권 100원을 기준으로 환급금을 총지급배율로 표시하면 이는 가장
가까운 $0.1$배로 반올림하는 규칙에 해당한다.\footnote{국가법령정보센터의
현행 「한국마사회법 시행령」 제10조(환급금) 및 제11조 제1항(단수의 처리)을
원문 대조했다(2025년 1월 31일 시행본). 제11조 제1항은 ``환급금에 10원 미만의
단수가 있을 때에는 5원 미만은 계산하지 아니하고, 5원 이상은 10원으로
계산한다''고 규정한다. 공식 조문은
\href{https://www.law.go.kr/LSW/lsInfoP.do?lsiSeq=268449}{국가법령정보센터}에서
확인할 수 있다.} 따라서 비상한 게시값 $D$에 대해 실제 총지급배율의 허용구간을
$[D-0.05,D+0.05]$로 두며, 앞의 $0.1$ 게시격자 전수감사는 실제 파일의 표시구조가
이 법정 단수처리와 양립하는지 별도로 확인한다. 이 근거 때문에 Panel B의
반올림 포락선은 절사 규칙을 임의로 가정한 민감도 구간이 아니라 법정 환급금
단수처리를 가격자료에 옮긴 측정구간이다.

실제로 삼쌍승 상한 조합이 있는 경주는 15,963개로, 19,284개 경주의 대부분이다.
삼쌍승과 비교 승식 모두 상한이 없는 clean point sample은 목표 승식별로 3,321개
경주다. 따라서 이 표본의 점추정과 나머지 15,963개 경주의 정보를 버리지 않기
위해 두 결과를 공동 주결과로 보고한다. Panel A는 3,321개 clean 경주의 점추정과
표본불확실성을 제시한다. Panel B는 clean 경주와 상한 경주를 합친 전체 19,284개
경주에서 검열·반올림 구간이 허용하는 TV·MAE와 기준모형 순위의 보수적
부분식별 경계를 제시한다. Panel A는 정확한 점비교를 제공하지만 전체 시장의
대표표본으로 간주하지 않으며, Panel B는 보조적 강건성 분석이 아니다. 네 목표
승식의 clean 표본은 같은 3,321개 race ID 집합임을 strict 감사로 확인한다. 두
표본의 출전두수, 연도, 경마장, 상한 조합 수와 상한 직전 비검열 배당의
꼬리분포를 먼저 비교하여 표본선택의 방향을 보고한다. 이 자료 감사는 아직 승식
간 재구성 성과를 계산하지 않으므로, 표 \ref{tab:data-audit}의 수치를 가격
정합성의 증거로 해석하지 않는다.